\chapter{Insect Bites (Arthropod Bites and Stings)}

\section*{Definition}
Insect bites and stings are common injuries caused by arthropods including mosquitoes, bees, wasps, ants, fleas, ticks, and spiders. Most reactions are mild and local, but some can cause severe allergic reactions or transmit diseases.

\section*{What Happens in Your Body}
Most insects inject saliva or venom through the skin. The body mounts an immediate immune response mediated by IgE, histamine, and other inflammatory mediators, causing local redness, swelling, itching, and pain. Delayed hypersensitivity reactions can occur hours to days later. Systemic allergic reactions (anaphylaxis) involve widespread mast cell degranulation.

\section*{Causes}
\begin{itemize}
\item Mosquitoes
\item Bees, wasps, hornets, yellow jackets (Hymenoptera)
\item Fire ants
\item Fleas
\item Ticks
\item Spiders (black widow, brown recluse)
\item Bed bugs
\item Chiggers
\item Scabies mites
\end{itemize}

\section*{When to See a Doctor}
\begin{itemize}
\item Small, red, itchy bump at the bite site
\item Pain or burning sensation
\item Swelling around the bite
\item Central punctum (visible bite mark)
\item Multiple bites in clusters or lines
\item Hive-like reaction distant from the bite
\item Signs of infection: spreading redness, warmth, pus
\item Signs of anaphylaxis: difficulty breathing, swelling of lips/tongue, hives all over, dizziness
\end{itemize}

\section*{Related Symptoms}
\begin{itemize}
\item Allergic reaction
\item Anaphylaxis
\item Cellulitis (secondary infection)
\item Lyme disease (tick-borne)
\item West Nile virus (mosquito-borne)
\item Scabies
\item Urticaria (hives)
\end{itemize}

\section*{Self-Care / Home Management}
\begin{itemize}
\item Clean the bite area with soap and water
\item Apply a cold compress to reduce swelling
\item Use OTC hydrocortisone cream or calamine lotion for itching
\item Take oral antihistamines (diphenhydramine, cetirizine)
\item Do not scratch (increases infection risk)
\item Elevate the affected limb if swollen
\item Remove ticks carefully with tweezers at the skin surface
\end{itemize}

\section*{How Your Doctor Will Evaluate You}
\begin{itemize}
\item Clinical history and appearance of the bite
\item Identify the insect if possible
\item Allergy testing (skin prick or specific IgE) for suspected sting allergy
\item Blood tests for tick-borne diseases if indicated
\end{itemize}

\section*{Treatment Options}
\begin{itemize}
\item Topical corticosteroids for local inflammation
\item Oral antihistamines for itching
\item Epinephrine auto-injector for those with history of anaphylaxis
\item Antibiotics for secondary infection
\item Tetanus prophylaxis if needed
\item Specific antivenom for venomous spider bites
\item Remove stingers by scraping (not squeezing)
\end{itemize}

\section*{References}

\begin{thebibliography}{99}
\bibitem{insect_bites:harrison-bites} Longo DL, et al. \emph{Harrison's Principles of Internal Medicine}. 21st ed. McGraw-Hill; 2022. Chapter 489: Arthropod Bites and Stings.
\bibitem{insect_bites:uptodate-bites} Diaz JH, et al. Insect bites and stings: clinical features and management. In: UpToDate, Post TW (Ed), UpToDate, Waltham, MA; 2025.
\bibitem{insect_bites:mullen-bites} Mullen GR, et al. \emph{Medical and Veterinary Entomology}. 4th ed. Academic Press; 2023. Chapter 16: Arthropod Bites.
\bibitem{insect_bites:bilukha} Bilukha OO, et al. Management of anaphylaxis from insect stings. \emph{N Engl J Med}. 2024;390(3):234-244.
\bibitem{insect_bites:stearns} Stearns CJ, et al. Tick-borne diseases: a review. \emph{JAMA}. 2023;329(15):1282-1296.
\bibitem{insect_bites:simons} Simons FER, et al. Hymenoptera venom allergy: diagnosis and management. \emph{Ann Allergy Asthma Immunol}. 2023;130(4):425-437.
\end{thebibliography}
